\section{几何与概率公式推导}
\subsection{凸多边形直径在顶点对上取得}
设 $X=\sum_i\lambda_iV_i$、$Y=\sum_j\mu_jV_j$，其中系数非负且分别和为1。则
\begin{equation}
\norm{X-Y}=\left\lVert\sum_{i,j}\lambda_i\mu_j(V_i-V_j)\right\rVert
\le\sum_{i,j}\lambda_i\mu_j\norm{V_i-V_j}
\le\max_{i,j}\norm{V_i-V_j}.
\end{equation}
反之，顶点属于多边形，最大顶点距离必能取到，故式\eqref{eq:diameter}成立。圆盘是凸集，所以全部顶点在圆内等价于整个凸多边形在圆内。

\subsection{最小包围圆半径与直径的关系}
包含一对距离为 $D$ 的点需要半径至少 $D/2$。最小包围圆的圆心处于支撑点的凸包内，否则可稍移圆心使所有支撑距离同时减少。在平面上，圆心可由至多三个支撑点的凸组合表示。

两支撑点时，它们为对径点，$r_*=D/2$。三支撑点时，形成包含圆心的三角形，其最大角 $\gamma$ 位于 $[60^\circ,90^\circ]$，由正弦定理最大边长为 $2r_*\sin\gamma\ge\sqrt3r_*$。由于该边不超过 $D$，故 $r_*\le D/\sqrt3$，等边三角形取等号。

三点 $A=(a_x,a_y),B=(b_x,b_y),C=(c_x,c_y)$ 的外接圆心 $U$ 可由
\begin{equation}
2\begin{pmatrix}b_x-a_x&b_y-a_y\\c_x-a_x&c_y-a_y\end{pmatrix}U
=\begin{pmatrix}\norm B^2-\norm A^2\\\norm C^2-\norm A^2\end{pmatrix}
\end{equation}
求解。行列式接近零时应改查最远点对，避免除以近零数。求得圆心后，认证半径统一重算为 $\max_i\norm{U-V_i}$。

\subsection{圆盘外包络的误差}
半径 $r$ 的圆，其外接正 $m$ 边形顶点半径为 $r\sec(\pi/m)$，最大径向外扩为
\begin{equation}
\Delta_m=r\bigl(\sec(\pi/m)-1\bigr)
\simeq\frac{r\pi^2}{2m^2}.
\end{equation}
例如 $r=1800,m=96$ 时约为 $0.964\,\mathrm m$。外包络包含真值，故对它认证清除是保守的；但接近20米阈值时可增加边数或局部精确处理圆弧。该径向外扩不能直接作为任意近乎平行交会后的最终直径误差界，交会几何可能放大误差。

\subsection{三圆交保证区域的证明}
平移旋转后令 $S_1=0,\theta_1=0$，候选点记为向量 $s$，$u_\pm=e(\pm\eps)$。若 $s\in\mathcal V_0$，则
\begin{equation}
\norm s\le a,\qquad \norm{s-au_\pm}\le a
\ \Longrightarrow\ s\cdot u_\pm\ge\frac{\norm s^2}{2a}.
\end{equation}
对 $|\phi|\le\eps$，有 $e(\phi)=\lambda u_-+\mu u_+$，其中
\begin{equation}
\lambda=\frac{\sin(\eps-\phi)}{\sin(2\eps)},\quad
\mu=\frac{\sin(\eps+\phi)}{\sin(2\eps)},\quad
\lambda+\mu=\frac{\cos\phi}{\cos\eps}\ge1.
\end{equation}
故 $s\cdot e(\phi)\ge\norm s^2/(2a)$，进而 $\norm{s-ae(\phi)}\le a$。当 $0\le r\le a$，$re(\phi)$ 是 $0$ 与 $ae(\phi)$ 的凸组合，由距离凸性得 $\norm{s-re(\phi)}\le a$。当 $a\le r\le b$，
\begin{equation}
\norm{s-re(\phi)}^2-r^2
=\norm s^2-2r\,s\cdot e(\phi)
\le\norm s^2(1-r/a)\le0.
\end{equation}
两段分别得到 $d_2\le a$ 与 $d_2\le r$，故总有 $d_2\le\max(a,r)\le\rho$。

完整扇形模型中，取可能源位置 $0,au_-,au_+$ 和半径 $a$，可分别推出三个圆盘约束的必要性。若严格排除 $r\le5$ 或裁去圆域外区域，则这些测试位置可能不再允许，所以此时只保留充分性结论。这个证明也避免了把圆弧扇形的极点误写成只有三个点。

\subsection{径向后验及第二次接收概率}
固定源位置 $G$ 时，允许的接收半径区间是 $[\max(a,d_1),b]$，其长度即 $K_1(d_1)$。两次接收共用半径，所以同时允许的区间长度为 $K_{12}(d_1,d_2)$。这证明式\eqref{eq:weights}，也说明为何不能将两个独立半径的权重相乘。

对圆心第一点，交换积分顺序得
\begin{align}
I_1&=\int_a^b\int_0^\rho r\,\mathrm dr\,\mathrm d\rho
=\frac{b^3-a^3}{6},\\
I_2&=\int_a^b\int_0^\rho r^2\,\mathrm dr\,\mathrm d\rho
=\frac{b^4-a^4}{12}.
\end{align}
严格排除5米近距区后，分母、分子分别变为 $I_1-(b-a)5^2/2$ 与 $I_2-(b-a)5^3/3$。

给定候选点 $S$ 和第二误差 $\delta\in[-\eps,\eps]$，令 $\theta_2=\arg(G-S)+\delta$。任意区域指标 $H$ 的第二次示向条件期望为
\begin{equation}
\mathbb E[H\mid\text{第二次示向}]=
\frac{\displaystyle\int\!\int rK_{12}\ind_{d_2>5}
\left[\frac1{2\eps}\int_{-\eps}^{\eps}H(K_2(S,\theta_2))\,\mathrm d\delta\right]\mathrm dr\,\mathrm d\phi}
{\displaystyle\int\!\int rK_{12}\ind_{d_2>5}\,\mathrm dr\,\mathrm d\phi}.
\end{equation}
严格按第一读数为示向时，径向积分从5开始。取 $H=r_*$ 得平均包围半径，取 $H=\ind_{r_*\le20}$ 得示向条件下的一次认证清除概率；近距分支另加其概率。

对一维积分 $\int_l^u f(t)\,\mathrm dt$，Gauss--Legendre求积写为
\begin{equation}
\int_l^u f(t)\,\mathrm dt\approx\frac{u-l}{2}\sum_{k=1}^q w_k
f\left(\frac{u+l}{2}+\frac{u-l}{2}\xi_k\right).
\end{equation}
将径向、角度和第二误差分别离散，形成张量求积。先在 $r=a$、$d_2=a,b,5$ 及圆域边界处划分区间，能减轻折点影响。确定性求积仍有离散误差，不应称为解析精确值；用加倍节点数比较收敛，并与独立抽样对照。

\subsection{局部交会面积与直径}
取两个单位法向夹角为 $\alpha$，误差带满足 $|n_1\cdot\Delta|\le h_1$、$|n_2\cdot\Delta|\le h_2$。线性映射 $\Delta\mapsto(n_1\cdot\Delta,n_2\cdot\Delta)$ 的行列式绝对值为 $|\sin\alpha|$，故面积为 $4h_1h_2/|\sin\alpha|$。

选 $n_1=(1,0)$、$n_2=(\cos\alpha,\sin\alpha)$，四个顶点为
\begin{equation}
\Delta_{s,t}=\left(s h_1,\frac{t h_2-s h_1\cos\alpha}{\sin\alpha}\right),\qquad s,t\in\{-1,1\}.
\end{equation}
区域中心对称，直径等于最大顶点范数的两倍；选择 $st$ 使交叉项为正即得式\eqref{eq:linear-diameter}。

\subsection{一般正多边形布局的覆盖半径}
设圆心加半径 $r_h\le R_0$ 的正 $n$ 边形顶点，$n\ge3$，令 $\beta=\pi/n$。把平面划为相邻顶点的对称扇区。半径为 $t$ 的点，到最近环顶点的最坏距离出现在扇区中线，其平方为
\begin{equation}
g(t)^2=t^2+r_h^2-2tr_h\cos\beta.
\end{equation}
同时圆心距离为 $t$。两者相等于 $t_0=r_h/(2\cos\beta)$。在 $[0,t_0]$ 上，最近扫描点距离不超过 $t_0$；在 $[t_0,R_0]$ 上，$g(t)^2$ 为凸二次函数，其最大值在端点取得。因此
\begin{equation}
\boxed{c_n(r_h)=\max\left\{\frac{r_h}{2\cos(\pi/n)},
\sqrt{R_0^2+r_h^2-2R_0r_h\cos(\pi/n)}\right\}.}
\label{eq:generalcover}
\end{equation}
该式适用于上述 $n\ge3,\ 0<r_h\le R_0$ 条件，内部项不是 $r_h\sin(\pi/n)$。取六边形、环半径1200米，即得基线覆盖半径968.90米；取八边形、环半径1010米，即得 matrix 布局的949.14米。

若要求 $c_n(r_h)\le a$，必须满足
\begin{equation}
R_0\cos\beta-\sqrt{a^2-R_0^2\sin^2\beta}\le r_h
\le R_0\cos\beta+\sqrt{a^2-R_0^2\sin^2\beta},
\end{equation}
同时满足 $r_h\le2a\cos\beta$ 和 $r_h\le R_0$；根号内为负时，该布局无解。六边形下最小可行环半径为
\begin{equation}
r_{h,\min}=900\sqrt3-100\sqrt{19}\approx1122.96\,\mathrm m.
\end{equation}
圆心起步、沿环访问的开放路径长度为 $r_h+2(n-1)r_h\sin(\pi/n)$。以此加上 $(n+1)$ 个停点的检测和切频代价，可以比较同族布局，但不构成所有搜索路径的全局最优性证明。

\subsection{栅格证书应覆盖整个单元}
若正方形单元边长为 $h$、中心为 $z$，则其中任意点到 $z$ 的距离不超过 $h/\sqrt2$。要证明整个单元位于以检测点 $q$ 为圆心的接收圆内，可以检查
\begin{equation}
\norm{q-z}+h/\sqrt2\le a.
\end{equation}
也可检查正方形四个顶点均在圆盘内，因为圆盘是凸集。只检查格心会把边缘未覆盖位置误判为已覆盖。

对联合半径剪枝，记单元 $B$ 到正观测点的最小距离为 $d_i^-(B)$、到负观测点的最大距离为 $d_k^+(B)$。若
\begin{equation}
\max\left(a,\max_i d_i^-(B)\right)>
\min\left(b,\min_k d_k^+(B)\right),
\end{equation}
则整格不存在可行半径，可安全删除；反向不成立，未被删去的格仍可能不可行，应细分而不是直接认定有源。

\subsection{知识矩阵的证据更新与单调性}
设真实位置与半径为 $(G_j,\rho_j)$。对时刻 $t$ 的反馈 $y_t$，按状态表生成约束集 $\mathcal H(y_t,p_t)$，联合知识更新为
\begin{equation}
\mathcal F_j^{t+1}=\mathcal F_j^t\cap\mathcal H(y_t,p_t).
\end{equation}
若反馈符合题设且 $\mathcal F_j^0$ 包含真值，每个新约束也包含真值，故归纳可知真值始终不被删除。与此同时 $\mathcal F_j^{t+1}\subseteq\mathcal F_j^t$，新信息不会扩大真实可行集。

这一性质要求保留全部有效历史。若一个格的最新证据覆盖掉另一真实位置的旧证据，重建集合可能重新扩大；若把光学失败错误转为无信号，甚至可能删除真值。因此紧凑矩阵应作为历史事件的索引和摘要，完整几何证据宜保留为追加记录。整行标记已清除是任务状态更新，不应抹掉用于复盘的历史内容。

\subsection{半径界的安全性与空域检查}
对任意正观测 $p$，因 $G_j\in\widehat K_j$，有
\begin{equation}
\min_{X\in\widehat K_j}\norm{p-X}\le\norm{p-G_j}\le\rho_j.
\end{equation}
对任意无线电负观测 $q$，有
\begin{equation}
\rho_j<\norm{q-G_j}\le\max_{X\in\widehat K_j}\norm{q-X}.
\end{equation}
仅在全向场景，分别取观测的最大下界和最小上界得到
\begin{equation}
L_j=\max\{a,\max_{p\in\mathcal P_j}\min_{X\in\widehat K_j}\norm{p-X}\},\quad
U_j=\min\{b,\min_{q\in\mathcal N_j}\max_{X\in\widehat K_j}\norm{q-X}\}.
\label{eq:radiusbounds}
\end{equation}
定向场景的无线电无信号不能给出上述上界。若算出 $L_j>U_j$，应检查证据、坐标、事件类别及数值误差，不能通过强行截断掩盖矛盾。

例如真实源在 $(30,0)$、接收半径1000米，在原点光学清除失败，但无线电仍可接收。光学失败只排除20米圆，若将它解释为1000米圆内无信号，会直接删去真实源。矩阵统一记录多类反馈，并不意味着各类反馈共享同一种距离约束。

\subsection{位掩码缺口与有限贪心补测}
对每个频道与每个格组成的任务对 $(j,\ell)$，定义待确认有限全集
\begin{equation}
\mathcal E=\{(j,\ell):\ell\in\mathcal U_j\}.
\end{equation}
候选停点 $S$ 可服务其中整格满足接收保证的任务对，记为 $\mathcal E(S)$。若候选集有限，且已经验证 $\mathcal E\subseteq\bigcup_S\mathcal E(S)$，每次选择新增任务数为正的点并实际完成对应频道测量，则未确认任务数至少减少1。因此至多 $|\mathcal E|$ 步，或至多遍历整个有效候选集合，就能消除当前发现缺口。

这是有限性结论，不是最短路径或最少停点定理；动态候选集和有限预算下仍应检查实际剩余缺口。若仅使用格心构造三角网格，则需额外验证整格条件，不能直接套用连续覆盖半径的结论。

在每次真实动作完成后，先更新矩阵，再更新候选格和覆盖格；只登记真正返回反馈的频道。由于预算中断而未执行的后续频道，不能提前加入已确认集合。这样位掩码的零缺口才具有发现层证书意义。

\subsection{首次示向的有限光学覆盖构造}
以首次示向为局部横轴，真源位于 $[0,1500]\times[-1500\sin\eps,1500\sin\eps]$ 中。将该矩形划为75列、3行，每格横向20米、纵向约17.45米，半对角线约13.27米，小于20米。依次访问225个格心，若真实源仍在该区域且无预算中断，必有一次清除成功。该构造说明有限兜底存在；最终程序实际采用当前可行域的25米相交方格，默认至多64格，两者不可混作同一实现。
